%% Generated by Sphinx.
\def\sphinxdocclass{report}
\documentclass[letterpaper,10pt,turkish]{sphinxmanual}
\ifdefined\pdfpxdimen
   \let\sphinxpxdimen\pdfpxdimen\else\newdimen\sphinxpxdimen
\fi \sphinxpxdimen=.75bp\relax
\ifdefined\pdfimageresolution
    \pdfimageresolution= \numexpr \dimexpr1in\relax/\sphinxpxdimen\relax
\fi
\newdimen\sphinxremdimen\sphinxremdimen = 10pt
%% let collapsible pdf bookmarks panel have high depth per default
\PassOptionsToPackage{bookmarksdepth=5}{hyperref}

\PassOptionsToPackage{booktabs}{sphinx}
\PassOptionsToPackage{colorrows}{sphinx}

\PassOptionsToPackage{warn}{textcomp}
\usepackage[utf8]{inputenc}
\ifdefined\DeclareUnicodeCharacter
% support both utf8 and utf8x syntaxes
  \ifdefined\DeclareUnicodeCharacterAsOptional
    \def\sphinxDUC#1{\DeclareUnicodeCharacter{"#1}}
  \else
    \let\sphinxDUC\DeclareUnicodeCharacter
  \fi
  \sphinxDUC{00A0}{\nobreakspace}
  \sphinxDUC{2500}{\sphinxunichar{2500}}
  \sphinxDUC{2502}{\sphinxunichar{2502}}
  \sphinxDUC{2514}{\sphinxunichar{2514}}
  \sphinxDUC{251C}{\sphinxunichar{251C}}
  \sphinxDUC{2572}{\textbackslash}
\fi
\usepackage{cmap}
\usepackage[T1]{fontenc}
\usepackage{amsmath,amssymb,amstext}
\usepackage{babel}



\usepackage{tgtermes}
\usepackage{tgheros}
\renewcommand{\ttdefault}{txtt}



\usepackage[Sonny]{fncychap}
\ChNameVar{\Large\normalfont\sffamily}
\ChTitleVar{\Large\normalfont\sffamily}
\usepackage{sphinx}

\fvset{fontsize=auto}
\usepackage{geometry}


% Include hyperref last.
\usepackage{hyperref}
% Fix anchor placement for figures with captions.
\usepackage{hypcap}% it must be loaded after hyperref.
% Set up styles of URL: it should be placed after hyperref.
\urlstyle{same}

\addto\captionsturkish{\renewcommand{\contentsname}{Contents:}}

\usepackage{sphinxmessages}
\setcounter{tocdepth}{1}



\title{Gömülü Linux Sistemleri \sphinxhyphen{} Geliştirme ve Uygulama}
\date{04 Haz 2026}
\release{1.0.0}
\author{Kaan Aslan}
\newcommand{\sphinxlogo}{\vbox{}}
\renewcommand{\releasename}{Yayım}
\makeindex
\begin{document}

\ifdefined\shorthandoff
  \ifnum\catcode`\=\string=\active\shorthandoff{=}\fi
  \ifnum\catcode`\"=\active\shorthandoff{"}\fi
\fi

\pagestyle{empty}
\sphinxmaketitle
\pagestyle{plain}
\sphinxtableofcontents
\pagestyle{normal}
\phantomsection\label{\detokenize{index::doc}}


\sphinxAtStartPar
Add your content using \sphinxcode{\sphinxupquote{reStructuredText}} syntax. See the
\sphinxhref{https://www.sphinx-doc.org/en/master/usage/restructuredtext/index.html}{reStructuredText}
documentation for details.

\sphinxstepscope


\chapter{Giriş}
\label{\detokenize{introduction:giris}}\label{\detokenize{introduction::doc}}
\sphinxAtStartPar
Bu bölümde temel kavramlar tanıtılacaktır.


\section{Gerekli Malzemeler}
\label{\detokenize{introduction:gerekli-malzemeler}}
\sphinxAtStartPar
Kitaptaki uygulamaları gerçekleştirebilmek için gerekli olabilecek malzemeler şunlardır:
\begin{itemize}
\item {} 
\sphinxAtStartPar
\sphinxstylestrong{Raspberry Pi SBC} (\sphinxstyleemphasis{Single Board Computer}): Version 3, 4 ya da 5 olabilir. Yeni satın alacakların
Raspberry Pi 5 almalarını tavsiye ediyoruz.

\item {} 
\sphinxAtStartPar
\sphinxstylestrong{Raspberry Pi 5 soğutucusu}: Raspberry Pi 5 alınacaksa soğutucu ya da soğutuculu kılıf da alınmalı.
Raspberry Pi 3 ve 4 için soğutucu alınmayabilir.

\item {} 
\sphinxAtStartPar
\sphinxstylestrong{Micro SD Kart}: Raspberry Pi için 64 GB yüksek hızlı micro SD kart tavsiye ediyoruz. 32 GB de
olabilir. Ancak 32 GB’den düşüğü tavsiye etmiyoruz.

\item {} 
\sphinxAtStartPar
\sphinxstylestrong{HDMI Kablosu}: Raspberry Pi 4 ve Raspberry Pi 5 için \sphinxstyleemphasis{Micro HDMI \sphinxhyphen{}\textgreater{} HDMI Kablosu}, Raspberry Pi 3
için \sphinxstyleemphasis{HDMI \sphinxhyphen{}\textgreater{} HDMI} Kablosu. Biz kursumuzda Raspberry Pi 5 kullanacağız.

\item {} 
\sphinxAtStartPar
\sphinxstylestrong{USB Kombo Klavye + Mouse}

\item {} 
\sphinxAtStartPar
\sphinxstylestrong{GPIO Breadboard Aktarma Kablosu}: Raspberry Pi için.

\item {} 
\sphinxAtStartPar
\sphinxstylestrong{Güç Kaynağı Adaptörü}: Raspberry Pi 3 ve 4 için 5V/3A, Raspberry Pi 5 için 5V/5A (27W).
Raspberry Pi 4 ve 5 USB Type C kullanıyor. Raspberry Pi 3 Micro USB kablosu kullanıyor.

\item {} 
\sphinxAtStartPar
\sphinxstylestrong{BeagleBone Black SBC} (\sphinxstyleemphasis{Single Board Computer}): Yeni satın alacak olanlar \sphinxstyleemphasis{BeagleBone Black 4G},
\sphinxstyleemphasis{BeagleBone Black Wireless} ya da \sphinxstyleemphasis{BeagleBone Black Industrial} alabilirler. Stokları kontrol etmek
gerekiyor. Biz kursumuzda \sphinxstyleemphasis{BeagleBone Black Industrial} kullanacağız.

\item {} 
\sphinxAtStartPar
\sphinxstylestrong{BeagleBone Black Güç Kaynağı}: Mini USB ile 5V/500mA güç adaptörü ya da yuvarlak girişli (DC Jack)
5V/1A güç adaptörü. Ancak BeagleBone Black USB kablo ile PC’den de güç alabilmektedir.

\item {} 
\sphinxAtStartPar
\sphinxstylestrong{BeagleBone Black Micro SD Kartı}: 16 GB olabilir.

\item {} 
\sphinxAtStartPar
\sphinxstylestrong{BeagleBone Black HDMI Dönüştürücü}: \sphinxstyleemphasis{Micro HDMI \sphinxhyphen{}\textgreater{} HDMI} dönüştürücü.

\item {} 
\sphinxAtStartPar
\sphinxstylestrong{Breadboard}: Standart boy breadboard ve isteğe bağlı küçük boy breadboard’lar da olabilir.

\item {} 
\sphinxAtStartPar
\sphinxstylestrong{Jumper Kablo Seti}

\item {} 
\sphinxAtStartPar
\sphinxstylestrong{USB \sphinxhyphen{}\textgreater{} UART Dönüştürücü}: CP2102 çevirici modül olabilir.

\item {} 
\sphinxAtStartPar
\sphinxstylestrong{DS3231 RTC Modülü}

\item {} 
\sphinxAtStartPar
\sphinxstylestrong{USB Uzatma Kablosu}: Gerekebilir.

\item {} 
\sphinxAtStartPar
\sphinxstylestrong{Ethernet RJ45 Kablosu}: Eğer SBC’nin wireless özelliği varsa buna gerek duymayabilirsiniz.

\end{itemize}


\section{Gömülü Sistemlere Giriş}
\label{\detokenize{introduction:gomulu-sistemlere-giris}}
\sphinxAtStartPar
Asıl amacı bilgisayar olmayan fakat bilgisayar devresi içeren sistemlere genel olarak
\sphinxstyleemphasis{gömülü sistemler} (\sphinxstyleemphasis{embedded systems}) denilmektedir. Yani gömülü sistemler başka amaçları
gerçekleştirmek için tasarlanmış olan aygıtların içerisindeki bilgisayar sistemleridir. Örneğin
elektronik tartılar, biyomedikal aygıtlar, GPS cihazları, turnike geçiş sistemleri (validatörler),
müzik kutuları, kapı güvenlik aygıtları, otomobiller içerisindeki kontrol panelleri birer gömülü
sistemdir. Gömülü sistemlerde en çok kullanılan programlama dili C’dir. Ancak genel amaçlı
işletim sistemlerinin yüklenebildiği SBC’lerde (\sphinxstyleemphasis{Single Board Computer}) diğer programlama dilleri
de kullanılabilmektedir.

\sphinxAtStartPar
Gömülü sistemler gömüldüğü aygıtta belli işlevleri sağlamak için kullanılmaktadır. Örneğin
buzdolaplarının, çamaşır makinelerinin içerisine yerleştirilen bilgisayar devreleri ve yazılımlar
onların önemli etkinliklerini yönetmekte, kullanıcı ile arayüz oluşturabilmektedir. Günümüzde
bilgisayar sistemi içeren aygıtlar artık her yeri sarmıştır. Bu nedenle gömülü sistemler üzerinde
geliştirme faaliyetleri de önemli bir iş alanı haline gelmiştir.


\subsection{Gömülü Sistemlerin Tipik Özellikleri}
\label{\detokenize{introduction:gomulu-sistemlerin-tipik-ozellikleri}}\begin{description}
\sphinxlineitem{\sphinxstylestrong{Belirli amaca yönelik işlemler:}}
\sphinxAtStartPar
Gömülü sistemler “genel değil belirli (\sphinxstyleemphasis{specific}) amaca yönelik” işlemleri gerçekleştirmektedir.
Yani genel amaçlı değil, özel amaçlı donanım ve yazılım hizmeti sunmaktadır. Bu sistemlerdeki
yazılımlar da genel amaçlı değil, belli bir amacı gerçekleştirmeyi hedeflemektedir.

\sphinxlineitem{\sphinxstylestrong{Düşük bilgi işlem kapasitesi:}}
\sphinxAtStartPar
Gömülü sistemler genellikle daha düşük bir bilgi işlem kapasitesine sahip bilgisayar
devreleridir. Örneğin bu sistemlerde kullanılan işlemciler genel amaçlı masaüstü işlemcilerden
genellikle daha yavaş olma eğilimindedir. Bu sistemlerdeki bellek miktarları (birincil ve ikinci
bellekler) genel amaçlı bilgisayar sistemlerine göre daha düşük olma eğilimindedir. Dolayısıyla
gömülü sistemlerin maliyetleri de genel amaçlı bilgisayar sistemlerine göre oldukça düşüktür.

\sphinxlineitem{\sphinxstylestrong{Düşük güç harcaması:}}
\sphinxAtStartPar
Gömülü sistemler genellikle (fakat her zaman değil) daha düşük güç harcamaktadır. Bu durum
onların bataryalarla beslenebilmesini dolayısıyla fiziksel taşınabilirliğini de artırmaktadır.
Tabii genel olarak gömülü sistemler düşük güç harcama eğiliminde olan sistemler olsalar da bu
durum her zaman böyle olmak zorunda değildir. Bazı gömülü sistemlerin gömüldüğü sistemlerde bir
güç kullanma sorunu yoktur. (Örneğin araba kantarı zaten bu işlevi gerçekleştirmek için önemli
bir güç harcamaktadır. Dolayısıyla bu sistemdeki gömülü sistemin harcadığı gücün önemi
olmayabilir.)

\sphinxlineitem{\sphinxstylestrong{Gerçek zamanlı olaylar:}}
\sphinxAtStartPar
Gömülü sistemlerin önemli bir bölümü \sphinxstyleemphasis{gerçek zamanlı} (\sphinxstyleemphasis{real time}) olaylarla ilişkilidir. Bu
sistemlerin belli bir bölümü dış dünyadaki değişimlere karşı bir yanıt oluşturmaya çalışmaktadır.
Örneğin bir gömülü sistem ortamdaki ısı belli bir kritik düzeye geldiğinde bir işlemi
başlatabilir. Ya da bir gömülü sistem kalp ritmi bozulduğunda kalbe uyarılar göndererek ritim
bozukluğunu düzeltmeye çalışabilir. Hava araçlarındaki gömülü sistemler o anki hava şartlarına
göre bir otomatik kumanda sistemi işlevini görüyor olabilir. Tabii gömülü sistemlerin bu gerçek
zamanlı işlem doğası bazı uygulamalarda \sphinxstyleemphasis{çok katı} (\sphinxstyleemphasis{hard realtime}) olurken bazı uygulamalarda
\sphinxstyleemphasis{daha gevşek} (\sphinxstyleemphasis{soft realtime}) olabilmektedir.

\sphinxlineitem{\sphinxstylestrong{Minimal kullanıcı arayüzü:}}
\sphinxAtStartPar
Gömülü sistemlerin bazılarında hiç girdi/çıktı birimi olmayabilir. Bazılarında ise girdi/çıktı
birimi olarak \sphinxstyleemphasis{düğmeler}, basit tuş takımları, küçük LCD’ler olabilir. Oysa genel amaçlı
bilgisayar sistemlerinde genellikle girdi/çıktı birimi olarak klavye, fare ve gelişmiş monitörler
kullanılmaktadır. Başka bir deyişle gömülü sistemler kullanıcı arayüzü bakımından minimal olma
eğilimindedir.

\sphinxlineitem{\sphinxstylestrong{Ucuz donanım:}}
\sphinxAtStartPar
Gömülü sistemlerdeki donanım birimleri nispeten ucuz olma eğilimindedir. Genel amaçlı
bilgisayarlara göre bunlar çok daha ucuz olarak temin edilebilmektedir.

\end{description}


\subsection{Gömülü Sistem Sınıflandırması}
\label{\detokenize{introduction:gomulu-sistem-siniflandirmasi}}
\sphinxAtStartPar
Gömülü sistemler \sphinxstyleemphasis{işlevsel gereksinimlere göre} ve \sphinxstyleemphasis{performans gereksinimlerine} göre de
sınıflandırılabilmektedir:

\begin{figure}[htbp]
\centering
\capstart

\noindent\sphinxincludegraphics[width=0.900\linewidth]{{embedded-classification}.png}
\caption{Gömülü sistem sınıflandırma şeması}\label{\detokenize{introduction:id1}}\end{figure}

\sphinxAtStartPar
\sphinxstylestrong{1) İşlevsel Gereksinimlere Göre Sınıflandırma:}
\begin{enumerate}
\sphinxsetlistlabels{\alph}{enumi}{enumii}{}{)}%
\item {} 
\sphinxAtStartPar
\sphinxstylestrong{Gerçek Zamanlı Olan ya da Gerçek Zamanlı Olmayan Gömülü Sistemler:} Bunlar \sphinxstyleemphasis{hard} ya da
\sphinxstyleemphasis{soft} real\sphinxhyphen{}time olabilmektedir.

\item {} 
\sphinxAtStartPar
\sphinxstylestrong{Bağımsız (Stand\sphinxhyphen{}alone) Olan ya da Bağımsız Olmayan Sistemler.}
Örneğin hesap makineleri, kapı güvenlik sistemleri, MP3 çalarlar gibi.

\item {} 
\sphinxAtStartPar
\sphinxstylestrong{Ağ (Network) Üzerinde İşlem Yapan Gömülü Sistemler:} Bunlar ağ işlemleri yapmak için
oluşturulmuş gömülü sistemlerdir. ATM makinelerini, ADSL Router gibi aygıtları bunlara
örnek verebiliriz.

\item {} 
\sphinxAtStartPar
\sphinxstylestrong{Mobil Aygıtlarda Kullanılan Gömülü Sistemler:} Bunlar küçük, taşınabilir aygıtlarda
kullanılan gömülü sistemlerdir. Telefonlar, GPS cihazları, dijital kameralar bunlara örnek
verilebilir.

\end{enumerate}

\sphinxAtStartPar
\sphinxstylestrong{2) Performans Gereksinimlerine Göre Sınıflandırma:}
\begin{enumerate}
\sphinxsetlistlabels{\alph}{enumi}{enumii}{}{)}%
\item {} 
\sphinxAtStartPar
Küçük Ölçekli (\sphinxstyleemphasis{Small Scale}) Gömülü Sistemler

\item {} 
\sphinxAtStartPar
Orta Ölçekli (\sphinxstyleemphasis{Medium Scale}) Gömülü Sistemler

\item {} 
\sphinxAtStartPar
Büyük Ölçekli (\sphinxstyleemphasis{Large Scale}) Gömülü Sistemler

\end{enumerate}

\sphinxAtStartPar
Biz burada \sphinxurl{https://www.ultralibrarian.com/2022/06/28/types-of-embedded-systems-characteristics-classifications-ulc}
adresindeki sınıflandırmayı kullandık. Ancak başka kaynaklarda başka türlü sınıflandırmalar da
yapılabilmektedir.


\subsection{Gömülü Sistemlerde İşlem Birimleri}
\label{\detokenize{introduction:gomulu-sistemlerde-islem-birimleri}}
\sphinxAtStartPar
Gömülü sistem yazılımlarının önemli bir bölümü bir işletim sistemi olmadan çalışacak biçimde
(\sphinxstyleemphasis{bare\sphinxhyphen{}metal}) olarak geliştirilmektedir. Bunun önemli nedenlerinden biri bunların kapasitelerinin
nispeten düşük olması diğeri de belirli bir amacı gerçekleştirecek biçimde tasarlanmış olmalarıdır.
Gömülü sistemlerin bir bölümünde \sphinxstyleemphasis{gerçek zamanlı işletim sistemleri} bir bölümünde ise \sphinxstyleemphasis{genel
amaçlı işletim sistemleri} kullanılmaktadır.

\sphinxAtStartPar
Gömülü sistemlerde genel olarak üç işlem birimi kullanılmaktadır:
\begin{enumerate}
\sphinxsetlistlabels{\arabic}{enumi}{enumii}{}{.}%
\item {} 
\sphinxAtStartPar
Mikrodenetleyiciler

\item {} 
\sphinxAtStartPar
Mikroişlemciler

\item {} 
\sphinxAtStartPar
DSP’ler

\end{enumerate}


\section{Mikroişlemci (CPU)}
\label{\detokenize{introduction:mikroislemci-cpu}}
\sphinxAtStartPar
Bir bilgisayar sisteminde aritmetik, mantıksal, bitsel işlemler ve karşılaştırma gibi bilgi işlem
faaliyetleri \sphinxstyleemphasis{mikroişlemci} (\sphinxstyleemphasis{microprocessor}) denilen birim tarafından yapılmaktadır.
Mikroişlemciler entegre devre biçiminde üretilmişlerdir. Mikroişlemcilere kavramsal olarak
\sphinxstyleemphasis{CPU} (\sphinxstyleemphasis{Central Processing Unit}) de denilmektedir. Aslında genel amaçlı bir bilgisayar sisteminde
komut çalıştıran pek çok işlemci bulunabilmektedir. CPU bu işlemcileri de programlayan \sphinxstyleemphasis{merkezi}
(\sphinxstyleemphasis{central}) işlemcidir. Bilgisayar sisteminde yerel bazı işlemlerden sorumlu yardımcı işlemciler de
vardır. Örneğin \sphinxstyleemphasis{kesme denetleyicisi} (\sphinxstyleemphasis{interrupt controller}), \sphinxstyleemphasis{disk denetleyicisi}
(\sphinxstyleemphasis{disk controller}), \sphinxstyleemphasis{DMA denetleyicisi} (\sphinxstyleemphasis{DMA controller}) gibi.


\section{Mikrodenetleyici (MCU)}
\label{\detokenize{introduction:mikrodenetleyici-mcu}}
\sphinxAtStartPar
Mikrodenetleyiciler tek bir chip içerisine yerleştirilmiş bir bilgisayar sistemi gibi düşünülebilirler.
Tipik olarak bir mikrodenetleyicide bir \sphinxstyleemphasis{işlemci} (\sphinxstyleemphasis{processor}), kendi içerisinde \sphinxstyleemphasis{RAM ve Flash EPROM},
dış dünya ile haberleşmek için kullanılan \sphinxstyleemphasis{Giriş/Çıkış (IO) birimleri} ve bazı \sphinxstyleemphasis{çevre birimleri}
(\sphinxstyleemphasis{peripherals}) bulunmaktadır. Mikrodenetleyicilere İngilizce \sphinxstyleemphasis{Microcontroller} ya da
\sphinxstyleemphasis{Microcontroller Unit (MCU)} da denilmektedir.

\begin{figure}[htbp]
\centering
\capstart

\noindent\sphinxincludegraphics[width=0.600\linewidth]{{microcontroller-internal}.png}
\caption{Tipik bir mikrodenetleyicinin iç yapısı}\label{\detokenize{introduction:id2}}\end{figure}

\sphinxAtStartPar
Mikrodenetleyicilerin işlem kapasiteleri ve içerdikleri bellek miktarları düşük olma eğilimindedir.
Ancak bunlar çok kolay programlanıp uygulamaya sokulabilmektedir. Mikrodenetleyicilere
\sphinxstyleemphasis{tek çiplik bilgisayar} (\sphinxstyleemphasis{single chip computer}) da denilmektedir. Mikrodenetleyiciler özellikle
gömülü sistemlerde tercih edilmektedir. Bunların düşük güç harcaması ve ucuz olmaları en büyük
avantajlarındandır. Gömülü uygulamalarda kullanılan pek çok mikrodenetleyici ailesi vardır.
Örneğin:
\begin{itemize}
\item {} 
\sphinxAtStartPar
Microchip PIC Mikrodenetleyici Ailesi (Microchip)

\item {} 
\sphinxAtStartPar
Renesas Mikrodenetleyici Ailesi (Renesas)

\item {} 
\sphinxAtStartPar
ARM Mikrodenetleyici Ailesi (Tasarımcısı ARM Holding, ancak çok çeşitli firmalar tarafından
üretiliyor)

\item {} 
\sphinxAtStartPar
AVR Mikrodenetleyici Ailesi (Atmel, ancak Atmel firması 2016’da Microchip tarafından satın alındı)

\item {} 
\sphinxAtStartPar
MSP Mikrodenetleyici Ailesi (Texas Instruments)

\end{itemize}


\section{SoC (System on Chip)}
\label{\detokenize{introduction:soc-system-on-chip}}
\sphinxAtStartPar
Bazı firmalar ayrı birimler olarak tasarlanmış mikroişlemcileri, RAM’leri, ROM’ları ve diğer bazı
çevre birimlerini tek bir entegre devrenin içerisine sıkıştırmaktadır. Bunlara genel olarak
\sphinxstyleemphasis{SoC} (\sphinxstyleemphasis{System on Chip}) denilmektedir. SoC mikrodenetleyicilere benzese de aslında onlardan
farklıdır. SoC’lar içerisindeki işlemcilerin ve belleklerin kapasiteleri yüksektir. Bunlar özel
amaçlı üretilirler ve bunların devrelerde kullanılmaları mikrodenetleyiciler kadar kolay değildir.
Bunların en önemli avantajları az yer kaplamasıdır. Örneğin Raspberry Pi kitlerinde Broadcom
isimli firmanın 2835, 2836, 2837, 2711, 2712 numaralı SoC çipleri kullanılmıştır. SoC’ların RAM
ve ROM bellek içermesi zorunlu değildir. Bazı SoC’lar RAM içerirken bazıları içermeyebilmektedir.
Örneğin Raspberry Pi 1, 2, 3 modellerinde kullanılan SoC’lar RAM içerirken Raspberry Pi 4 ve 5
modellerinde kullanılan SoC’lar RAM içermemektedir.


\section{SoM (System on Module)}
\label{\detokenize{introduction:som-system-on-module}}
\sphinxAtStartPar
SoC kavramına benzer diğer bir kavram da \sphinxstyleemphasis{SoM} (\sphinxstyleemphasis{System on Module}) kavramıdır. SoM bir işlemci
ve onunla ilişkili bazı birimlerin monte edildiği kartları belirtmek için kullanılmaktadır. SoM’lar
SoC’ları içerebilir. Ancak başka birimleri de içerebilir. Örneğin bir SoM bir işlemci, buna ilişkin
RAM, IO denetleyicisi (\sphinxstyleemphasis{IO controllers}) içeren bir kart olabilir. Örneğin \sphinxstyleemphasis{Raspberry Pi Pico} ve
\sphinxstyleemphasis{Raspberry Pi Compute Module} birer SBC’den ziyade birer SoM olarak ele alınabilir. SoM kavramını
zihninizde SoC ile SBC arasında bir yerde konumlandırabilirsiniz. SoM terimine yakın başka bir
kavram da \sphinxstyleemphasis{modül} (\sphinxstyleemphasis{module}) kavramıdır. Modül bir kart üzerine yerleştirilmiş çeşitli
bileşenlerden oluşan devrelere denilmektedir. Örneğin bir RTC kartına RTC modülü ya da bir GPS
kartına GPS modülü de denilmektedir.


\section{SBC (Single Board Computer)}
\label{\detokenize{introduction:sbc-single-board-computer}}
\sphinxAtStartPar
Küçük bir kit (baskılı devre) üzerine monte edilmiş bilgisayarlara \sphinxstyleemphasis{SBC} (\sphinxstyleemphasis{Single Board Computer})
denilmektedir. Genellikle bu kitlerde SoC’lar, RAM’ler, başka çevre birimleri ve IO işlemleri için
uçlar bulunmaktadır. Örneğin Raspberry Pi’lar, Beagleboard’lar SBC’lere örnek verilebilir.
SBC’ler klavye, fare ve monitör takılarak bir masaüstü bilgisayar gibi kullanılabilmektedir.
SBC’ler masaüstü bilgisayarlar gibi de kullanılabildiğinden bunlara Linux başta olmak üzere,
Android ve Windows gibi işletim sistemleri yüklenebilmektedir.


\section{CISC ve RISC İşlemci Mimarileri}
\label{\detokenize{introduction:cisc-ve-risc-islemci-mimarileri}}
\sphinxAtStartPar
Mikroişlemcileri tasarım mimarilerine göre \sphinxstyleemphasis{CISC} (\sphinxstyleemphasis{Complex Instruction Set Computers}) ve
\sphinxstyleemphasis{RISC} (\sphinxstyleemphasis{Reduced Instruction Set Computers}) olmak üzere iki kısma ayırabiliriz. CISC ailesi
işlemcilere Intel firmasının x86 ailesi işlemcilerini örnek verebiliriz. ARM, MIPS, PowerPC,
Itanium gibi işlemciler ise tipik RISC işlemcileridir. Her ne kadar CISC ve RISC isimleri komut
kümesi ile ilgili biçimde uydurulmuşsa da CISC ve RISC tasarımları başka bakımlardan da
farklılık göstermektedir.

\sphinxAtStartPar
Önceleri işlemcilerin (makro ve mikro) fazla sayıda makine komutuna sahip olması bir avantaj
olarak görülüyordu. Ancak daha sonraları bunun avantajdan ziyade dezavantaj oluşturduğu anlaşıldı.
Belli bir süreden sonra artık RISC mimarisinin CISC mimarisinden toplamda daha iyi bir tasarım
olduğu kabul görmüştür. Bu nedenle son dönem mikroişlemci tasarımlarında hep RISC mimarisi
kullanılmıştır.

\begin{figure}[htbp]
\centering
\capstart

\noindent\sphinxincludegraphics[width=0.600\linewidth]{{cisc-vs-risc}.png}
\caption{CISC ve RISC mimari karşılaştırması}\label{\detokenize{introduction:id3}}\end{figure}


\subsection{CISC ve RISC Arasındaki Temel Farklılıklar}
\label{\detokenize{introduction:cisc-ve-risc-arasindaki-temel-farkliliklar}}\begin{description}
\sphinxlineitem{\sphinxstylestrong{1) Komut sayısı ve karmaşıklığı:}}
\sphinxAtStartPar
CISC işlemcilerinde fazla sayıda makine komutu bulunmaktadır. Bu işlemcilerde bazı komutlar
temel işlemleri yaparken bazıları karmaşık işlemler yapmaktadır. Halbuki RISC işlemcilerinde
az sayıda temel makine komutları bulunmaktadır. Bu makine komutları da daha fazla transistör
kullanılarak daha hızlı çalışacak biçime getirilmiştir. Dolayısıyla CISC işlemcilerindeki bazı
komutlar RISC işlemcilerinde birkaç komutla karşılanabilmektedir. Ancak bu durum sanıldığının
tersine daha hızlı bir çalışma sağlama potansiyelindedir.

\sphinxlineitem{\sphinxstylestrong{2) Yazmaç (register) sayısı:}}
\sphinxAtStartPar
CISC işlemcilerinde az sayıda yazmaç (\sphinxstyleemphasis{register}), RISC işlemcilerinde ise fazla sayıda yazmaç
bulunma eğilimindedir. Yazmaç sayıları az olunca yazmaçların tekrar tekrar aynı değerlerle
yüklenmesi gerekebilmektedir. Bu da derleyicinin nesneleri daha kısa sürede yazmaçlarda
tutmasına yol açmaktadır. Yine CISC işlemcilerindeki bazı komutlar ancak bazı özel yazmaçlarla
kullanılmaktadır. (Örneğin Intel x86 işlemcilerinde MUL ve DIV komutlarının bir operandı EAX
ya da RAX yazmacında bulunmak zorundadır.) Oysa RISC işlemcilerinde her işlem her yazmaçla
yapılabilmektedir.

\sphinxlineitem{\sphinxstylestrong{3) Komut uzunlukları:}}
\sphinxAtStartPar
CISC işlemcilerinde komutlar farklı uzunluklarda olabilmektedir. Örneğin Intel’in x86
ailesinde 1 byte olan makine komutları da vardır, 5 byte olan makine komutları da vardır,
11 byte olan hatta 15 byte olan makine komutları da vardır. Halbuki RISC işlemcilerinde genel
olarak tüm makine komutları eşit uzunluktadır. Örneğin ARM işlemcilerinde her makine komutu
4 byte uzunluktadır. Böylece işlemci komutları bellekten daha etkin bir biçimde çekip (fetch
işlemi) onları daha çabuk anlamlandırmaktadır. Belli bir sayı sonraki ya da önceki makine
komutlarının yerini belirleyebilmektedir.

\sphinxlineitem{\sphinxstylestrong{4) Pipeline mekanizması:}}
\sphinxAtStartPar
RISC işlemcilerinde pipeline mekanizması CISC işlemcilerine göre daha iyi gerçekleştirilmektedir.
Pipeline işlemcinin bir makine komutunu çalıştırırken sonraki komutlar üzerinde hazırlık
işlemlerini yapması anlamına gelmektedir. RISC tasarımı pipeline mekanizmasının daha iyi
yürütülmesine olanak sağlamaktadır.

\sphinxlineitem{\sphinxstylestrong{5) Load/Store mimarisi:}}
\sphinxAtStartPar
RISC işlemcileri load/store tarzı makine komutları kullanmaktadır. Bu işlemcilerde belleğe
erişim yapan makine komutlarıyla aritmetik, mantıksal ve bitsel işlem yapan makine komutları
birbirinden ayrılmıştır. Örneğin RISC işlemcilerinde ADD, SUB gibi makine komutlarının her
iki operandı da yazmaç olmak zorundadır. Halbuki CISC işlemcilerinde bu tür makine komutlarının
bir operandı yazmaç bir operandı bellek olabilmektedir. Örneğin:

\begin{sphinxVerbatim}[commandchars=\\\{\}]
\PYG{n}{a} \PYG{o}{=} \PYG{n}{b} \PYG{o}{+} \PYG{n}{c}\PYG{p}{;}
\end{sphinxVerbatim}

\sphinxAtStartPar
gibi bir işlem CISC işlemcisi ile şu makine komutlarıyla yapılabilmektedir:
\begin{itemize}
\item {} 
\sphinxAtStartPar
b’yi yazmaca çeken makine komutu

\item {} 
\sphinxAtStartPar
Yazmaçtaki b ile bellekteki c’yi toplayan makine komutu

\item {} 
\sphinxAtStartPar
Yazmaçtaki toplamı a’ya yerleştiren makine komutu

\end{itemize}

\sphinxAtStartPar
Oysa aynı işlem RISC işlemcilerinde şöyle gerçekleştirilmektedir:
\begin{itemize}
\item {} 
\sphinxAtStartPar
b’yi yazmaca çeken makine komutu

\item {} 
\sphinxAtStartPar
c’yi yazmaca çeken makine komutu

\item {} 
\sphinxAtStartPar
Yazmaçlardaki b ile c’yi toplayan makine komutu

\item {} 
\sphinxAtStartPar
Yazmaçtaki sonucu a’ya yerleştiren makine komutu

\end{itemize}

\sphinxAtStartPar
Belleğe erişim komutlarıyla diğer komutların birbirinden ayrıldığı işlemcilere \sphinxstyleemphasis{load/store}
işlemciler ya da \sphinxstyleemphasis{register\sphinxhyphen{}register} işlemciler de denilmektedir.

\sphinxlineitem{\sphinxstylestrong{6) Operand sayısı:}}
\sphinxAtStartPar
RISC işlemcileri genel olarak (ancak hepsi değil) üç operandlı makine komutlarını kullanmaktadır.
Oysa CISC işlemcileri genellikle iki operandlı makine komutlarını kullanır. İki operandlı
makine komutlarında işlemin sonucu operand olan yazmaçlardan birine yerleştirildiği için o
yazmaç bozulmaktadır. Böylece derleyici o yazmaçtaki değeri yeniden kullanmak istediğinde
onu yeniden yüklemek zorunda kalmaktadır. Örneğin \sphinxcode{\sphinxupquote{a = b + c}} işlemi 32 bit Intel
işlemcilerinde tipik olarak şöyle gerçekleştirilmektedir:

\begin{sphinxVerbatim}[commandchars=\\\{\}]
\PYG{n}{MOV} \PYG{n}{EAX}\PYG{p}{,} \PYG{o}{\PYGZlt{}}\PYG{n}{a}\PYG{l+s+s1}{\PYGZsq{}}\PYG{l+s+s1}{nın adresi\PYGZgt{}}
\PYG{n}{ADD} \PYG{n}{EAX}\PYG{p}{,} \PYG{o}{\PYGZlt{}}\PYG{l+s+sa}{b}\PYG{l+s+s1}{\PYGZsq{}}\PYG{l+s+s1}{nin adresi\PYGZgt{}}
\PYG{n}{MOV} \PYG{o}{\PYGZlt{}}\PYG{n}{c}\PYG{l+s+s1}{\PYGZsq{}}\PYG{l+s+s1}{nin adresi\PYGZgt{}, EAX}
\end{sphinxVerbatim}

\sphinxAtStartPar
Burada EAX işlemcinin bir yazmacıdır. ADD makine komutu bu yazmaçtaki değer ile bellekteki
b’yi toplamakta ve sonucu yine EAX yazmacına yerleştirmektedir. Dolayısıyla EAX yazmacındaki
a değeri artık kaybolacaktır. Bu değer yeniden kullanılmak istendiğinde ise yeniden yüklemenin
yapılması gerekecektir. Şimdi aynı işlemi ARM işlemcilerinde yapacak olalım:

\begin{sphinxVerbatim}[commandchars=\\\{\}]
\PYG{n}{LDR} \PYG{n}{R0}\PYG{p}{,} \PYG{o}{\PYGZlt{}}\PYG{n}{a}\PYG{l+s+s1}{\PYGZsq{}}\PYG{l+s+s1}{nın bellek adresi\PYGZgt{}}
\PYG{n}{LDR} \PYG{n}{R1}\PYG{p}{,} \PYG{o}{\PYGZlt{}}\PYG{l+s+sa}{b}\PYG{l+s+s1}{\PYGZsq{}}\PYG{l+s+s1}{nin bellek adresi\PYGZgt{}}
\PYG{n}{ADD} \PYG{n}{R2}\PYG{p}{,} \PYG{n}{R1}\PYG{p}{,} \PYG{n}{R0}
\PYG{n}{STR} \PYG{n}{R2}\PYG{p}{,} \PYG{o}{\PYGZlt{}}\PYG{n}{c}\PYG{l+s+s1}{\PYGZsq{}}\PYG{l+s+s1}{nin bellek adresi\PYGZgt{}}
\end{sphinxVerbatim}

\sphinxAtStartPar
Burada R0, R1 ve R2 işlemcinin yazmaçlarıdır. Görüldüğü gibi ARM işlemcilerinde load/store
komutları dışındaki komutlar üç operandlıdır. Bu da yazmaçlardaki değerlerin gerektiğinde
bozulmamasına yol açmaktadır.

\sphinxlineitem{\sphinxstylestrong{7) Güç tüketimi:}}
\sphinxAtStartPar
RISC işlemcileri yukarıda belirttiğimiz tasarım prensiplerinden dolayı toplamda daha az güç
harcama eğilimindedir. Bu da onların mobil aygıtlarda kullanılmasının önemli bir gerekçesini
oluşturmaktadır.

\sphinxlineitem{\sphinxstylestrong{8) Komut çalışma süreleri:}}
\sphinxAtStartPar
RISC işlemcilerinde makine komutlarının çalışma süreleri birbirine yakındır. Ancak CISC
işlemcilerinde makine komutlarının çalışma süreleri arasında önemli farklılıklar olabilmektedir.
Örneğin ARM işlemcilerinde toplama, çıkartma, çarpma gibi makine komutları 1 cycle’da
yapılmaktadır. Ancak load/store komutlarının, jump komutlarının ve özel bazı komutların çalışma
süreleri 1 cycle’dan fazla olabilmektedir. Oysa örneğin Intel x86 işlemcilerinde komut süreleri
değişik faktörlere bağlı olarak birbirinden çok farklılaşmaktadır.

\end{description}

\sphinxAtStartPar
RISC ve CISC mimarilerini bir spektrum olarak düşünmek gerekir. Örneğin MIPS işlemcileri ARM
işlemcilerine göre bu spektrumun daha fazla RISC tarafındadır. Intel’in x86 işlemcilerini kategori
olarak CISC grubu işlemciler olarak ele alsak da Pentium işlemcileri ile birlikte bu işlemcilerde
de RISC prensipleri gittikçe daha fazla kullanılır hale gelmiştir.


\section{ARM İşlemcileri}
\label{\detokenize{introduction:arm-islemcileri}}

\subsection{ARM’ın Tarihçesi}
\label{\detokenize{introduction:arm-in-tarihcesi}}
\sphinxAtStartPar
ARM işlemcilerinin tarihi \sphinxstyleemphasis{Acorn Computer} isimli İngiliz firmasına dayanmaktadır. Bu firma
80’li yılların başlarında \sphinxstyleemphasis{BBC Micro} isimli 64K’lık ev bilgisayarlarını yapmıştı. Bu bilgisayarlarda
Rockwell’in 8 bitlik 6502 işlemcileri kullanılıyordu. Firma daha sonra \sphinxstyleemphasis{Berkeley RISC} projesinden
etkilenerek kendi RISC işlemcilerini yapmaya odaklandı. Böylece ilk ARM modelleri tasarlanmış oldu.
Şirket 1990’da \sphinxstyleemphasis{Apple} ve \sphinxstyleemphasis{VLSI Technology} şirketleriyle ortaklıklar da kurarak ARM ismini aldı.
(Eskiden ARM \sphinxstyleemphasis{Acorn RISC Machine} isminden kısaltılıyordu. Fakat daha sonra bu firma kurulunca bu
kısaltma \sphinxstyleemphasis{Advanced RISC Machine} haline getirildi.) Apple firması bu yeni firmaya maddi destek
sağlamıştır. \sphinxstyleemphasis{VLSI Technology} firması ise ekipman tedarik etmiştir. Acorn ise az sayıda tasarım
mühendisini bu yeni firmaya aktarmıştır. 2016 yılında \sphinxstyleemphasis{SoftBank Group} ARM hisselerinin önemli bir
bölümünü aldı. 2018’de ARM’ın Çin şubesinin yarısından fazlasını \sphinxstyleemphasis{Chine Investment} şirketine sattı.
2020’de \sphinxstyleemphasis{NVidia}, ARM’ı SoftBank Group’tan satın almak istediyse de satış gerçekleşmedi. Bugün
\sphinxstyleemphasis{SoftBank Group} ARM’ın \%90 civarındaki hisselerine sahiptir. Geri kalan hisseleri kurucu
ortaklarda ve halka arzdadır.


\subsection{ARM Lisans Türleri}
\label{\detokenize{introduction:arm-lisans-turleri}}
\sphinxAtStartPar
ARM (İngilizce’de de \sphinxstyleemphasis{ey ar em} biçiminde değil \sphinxstyleemphasis{arm} biçiminde okuyorlar) bir tasarım firmasıdır.
Yani fabrikalara sahip değildir. ARM yaptığı tasarımları lisanslandırarak üretici firmalara
satmaktadır. ARM’ın uyguladığı dört tür lisans vardır:
\begin{description}
\sphinxlineitem{\sphinxstylestrong{Entegre Devre Lisansı} (\sphinxstyleemphasis{Full Custom License}):}
\sphinxAtStartPar
Bu lisans müşterinin mikrodenetleyici tasarımını kendi özel ihtiyaçlarına göre özelleştirmesine
ve optimize etmesine olanak tanır. Örneğin Apple gibi, Qualcomm gibi, Samsung gibi firmalar bu
lisansa sahiptir.

\sphinxlineitem{\sphinxstylestrong{Mimari Lisans} (\sphinxstyleemphasis{Architecture License}):}
\sphinxAtStartPar
Bu lisans ARM’nin genel mikroişlemci mimarisine erişim sağlayan lisanstır. Ancak bu lisans
müşterinin mikrodenetleyiciyi özelleştirmesine izin vermez. Yani bu lisansa sahip olanlar
işlemci tasarımını kullanarak üretim yapabilirler ancak onu özelleştiremezler.

\sphinxlineitem{\sphinxstylestrong{Çekirdek Lisansı} (\sphinxstyleemphasis{Processor IP License}):}
\sphinxAtStartPar
Bu lisans ile müşteri yalnızca belli bir ARM işlemcisini (\sphinxstyleemphasis{core}) üretebilmektedir.

\sphinxlineitem{\sphinxstylestrong{Geliştirici Lisansı} (\sphinxstyleemphasis{Development License}):}
\sphinxAtStartPar
Bu lisans ARM’nin tasarım araçlarına erişim sağlayan bir lisanstır. Müşteri bu araçları
kullanarak kendi işlemcilerini tasarlayabilir.

\end{description}


\subsection{ARM Terminolojisi: Core, Cortex ve Versiyon}
\label{\detokenize{introduction:arm-terminolojisi-core-cortex-ve-versiyon}}
\sphinxAtStartPar
ARM dünyasında çalışmak için bazı terimler hakkında bilgi sahibi olmak gerekir. ARM firması
kendine özgü terimler uydurmuştur. Bu dünyada en çok karşılaşılan terimler \sphinxstyleemphasis{core}, \sphinxstyleemphasis{cortex} ve
\sphinxstyleemphasis{version} terimleridir.

\sphinxAtStartPar
ARM dünyasında \sphinxstyleemphasis{core} terimi belli bir mikroişlemci tasarımını belirtmektedir. Bu tasarım üretici
firmalar tarafından fiziksel hale getirilmektedir. Bir grup \sphinxstyleemphasis{core} bir araya getirildiğinde ve
işlemciyle ilişkili birtakım birimler de bunlara eklendiğinde \sphinxstyleemphasis{cortex}’ler oluşmaktadır. ARM
firması ARM’ın klasik versiyonlarında \sphinxstyleemphasis{cortex} terimini kullanmamıştır. Cortex terimi ARM11’den
sonra kullanılmaya başlanmıştır. Bir cortex belli sayıda \sphinxstyleemphasis{core} içerebilir. Bu \sphinxstyleemphasis{core}’lar noktalı
sayılara (\sphinxstyleemphasis{floating point}) ilişkin işlem yapan birimlere sahip olabilirler. Cortex içerisindeki
core’ların belli hızları vardır. Core ve cortex terimleri daha çok çip düzeyindeki mimariyi
belirtmektedir. Cortex’lerdeki işlemcilerin (core’ların) sistem programcısını ilgilendiren
\sphinxstyleemphasis{komut kümeleri} (\sphinxstyleemphasis{instruction sets}) de vardır. Böylece bir grup cortex belli bir komut kümesi
ile kullanılabilecek biçimde tasarlanmıştır. Komut kümesi mimarisine \sphinxstyleemphasis{ARM Versiyonları} ya da
İngilizcesiyle \sphinxstyleemphasis{ARM ISA} (\sphinxstyleemphasis{Instruction Set Architecture}) denilmektedir. Biz sistem programcısı
olarak elimizdeki cortex’in hangi komut kümesini kullanan ARM versiyonuna ilişkin olduğunu
bilmek durumundayız. Burada dikkat edilmesi gereken nokta \sphinxstyleemphasis{Cortex teriminin donanımsal mimariyi,
versiyonun ise yazılımsal mimariyi belirtiyor} olmasıdır. O halde gömülü sistem geliştiricisi
olarak bizim ilgilendiğimiz ARM içeren kart ile ilgili şu iki özelliği biliyor olmamız gerekir:
\begin{enumerate}
\sphinxsetlistlabels{\arabic}{enumi}{enumii}{}{.}%
\item {} 
\sphinxAtStartPar
Kartımızda ARM’ın hangi cortex’i kullanılmıştır?

\item {} 
\sphinxAtStartPar
Bu cortex’in ilişkin olduğu versiyon numarası nedir?

\end{enumerate}

\sphinxAtStartPar
Yukarıda da belirttiğimiz gibi farklı cortex’ler aynı versiyon numarasını kullanabilmektedir.


\subsection{ARM Mimari Profilleri}
\label{\detokenize{introduction:arm-mimari-profilleri}}
\sphinxAtStartPar
ARM dünyasında üç mimari profili (\sphinxstyleemphasis{profile}) bulunmaktadır:

\begin{figure}[htbp]
\centering
\capstart

\noindent\sphinxincludegraphics[width=0.600\linewidth]{{arm-profiles}.png}
\caption{ARM mimari profilleri}\label{\detokenize{introduction:id4}}\end{figure}
\begin{enumerate}
\sphinxsetlistlabels{\arabic}{enumi}{enumii}{}{.}%
\item {} 
\sphinxAtStartPar
\sphinxstylestrong{A (Application) Profili}

\item {} 
\sphinxAtStartPar
\sphinxstylestrong{R (Realtime) Profili}

\item {} 
\sphinxAtStartPar
\sphinxstylestrong{M (Microcontroller) Profili}

\end{enumerate}

\sphinxAtStartPar
Bu profiller ilgili cortex’lerin hangi tür uygulamalarda ideal olarak kullanılabileceğini
belirtmektedir. Profil isimleri cortex’lerde ve versiyon numaralarında \sphinxcode{\sphinxupquote{\sphinxhyphen{}}} karakterinden sonra
A, R, M harfleriyle belirtilmektedir. Örneğin \sphinxstyleemphasis{Cortex\sphinxhyphen{}A8} ve \sphinxstyleemphasis{ARMv7\sphinxhyphen{}A} gibi. A profilleri
masaüstü işletim sistemlerinin çalıştırılabileceği, tamamen kişisel bilgisayar olarak
kullanılabilecek cortex’leri belirtmektedir. Biz kursumuzda gömülü Linux programcısı olarak bu
\sphinxstyleemphasis{A} profilleriyle çalışacağız. R profilleri \sphinxstyleemphasis{gerçek zamanlı} (\sphinxstyleemphasis{realtime}) uygulamalar için daha
uygun olabilecek cortex’leri belirtmektedir. R profilleri A profillerine göre çok daha seyrek
kullanılmaktadır. M profilleri ARM’ın mikrodenetleyici olarak kullanılan cortex’lerini
belirtmektedir. Bu cortex’ler Linux işletim sisteminin yüklenmesine izin vermemektedir. Bu
cortex’ler genellikle \sphinxstyleemphasis{bare\sphinxhyphen{}metal} programlarla ya da gerçek zamanlı işletim sistemleriyle
kullanılmaktadır.

\sphinxAtStartPar
ARM’ın iki önemli versiyonu ARMv7\sphinxhyphen{}A ve ARMv8\sphinxhyphen{}A’dır. ARMv7’ler 32 bitlik bir arayüz sunmaktadır.
Dolayısıyla bu komut mimarisini kullanan cortex’ler 32 bitlik işlemciler içermektedir. ARMv8\sphinxhyphen{}A
versiyonları ise 64 bitlik arayüz sunmaktadır. Dolayısıyla bu versiyonları kullanan cortex’ler
64 bitlik işlemcilere ilişkindir. Ancak ARM cortex’lerinin bir grubu hem 32 bit hem de 64 bit
olarak kullanılabilmektedir. Bu cortex’lerdeki işlemcileri kullanan işletim sistemleri 32 bitlik
ve 64 bitlik programları zaman paylaşımlı olarak bir arada çalıştırabilmektedir.


\section{İşlemcilerin Bit Genişliği}
\label{\detokenize{introduction:islemcilerin-bit-genisligi}}
\sphinxAtStartPar
Mikroişlemciler ilk çıktığında 8 bit işlemleri yapabiliyordu (örneğin 8080, 6800 gibi işlemciler).
Bunlara 8 bitlik mikroişlemci deniyordu. Sonra 16 bitlik mikroişlemciler çıktı (örneğin Intel’in
8086, 8088 işlemcileri gibi). Bunu 32 bit işlemciler izledi. Günümüzde artık ağırlıklı olarak 64
bit işlemciler kullanılmaktadır. Peki bir işlemcinin N bitlik olmasının anlamı nedir?
\begin{itemize}
\item {} 
\sphinxAtStartPar
\sphinxstylestrong{Tek hamlede işlem kapasitesi:} N bitlik bir işlemcide tek hamlede (yani tek bir makine
komutuyla) N bit üzerinde işlem yapılabilmektedir. Örneğin 32 bitlik bir mikroişlemcide tek
bir makine komutu ile 32 bitlik iki sayıyı toplayıp çarpabiliriz.

\item {} 
\sphinxAtStartPar
\sphinxstylestrong{Fiziksel RAM adresleme kapasitesi:} N bitlik bir mikroişlemci genellikle (ama her zaman
değil) 2$^{\text{N}}$ uzunlukta bir fiziksel RAM’i adresleyebilmektedir. Örneğin 32 bitlik bir
mikroişlemciye biz tipik olarak 2$^{\text{32}}$ = 4 GB’lik bir RAM bağlayabiliriz. Elimizdeki RAM
64 GB olsa bile 32 bit bir mikroişlemci bu RAM’in ancak ilk 4 GB’sini kullanabilmektedir. Tabii
bu durum her zaman böyle değildir. Örneğin 8086 işlemcisi 16 bit olduğu halde 2$^{\text{16}}$
değil 2$^{\text{20}}$ uzunluğunda (1 MB) fiziksel belleği adresleyebiliyordu. Benzer biçimde 64
bitlik Intel ve AMD işlemcileri 2$^{\text{64}}$ değil 2$^{\text{48}}$ uzunluğundaki RAM’leri
adresleyebilmektedir.

\item {} 
\sphinxAtStartPar
\sphinxstylestrong{CPU\sphinxhyphen{}RAM arası transfer genişliği:} N bitlik bir mikroişlemcide genellikle işlemci ile RAM
arasında transfer N bit olarak yapılmaktadır. Örneğin 32 bitlik bir mikroişlemcide tek bir
makine komutuyla RAM’den CPU’ya 32 bitlik bir veri transfer edilebilir. Ancak bu durum her
zaman böyle değildir. Örneğin 16 bitlik Intel 8086 işlemcisi bellek transferlerini 16 bit
olarak yaparken 8088 işlemcisi 8 bit yapıyordu.

\end{itemize}

\sphinxAtStartPar
8 bitten 16 bite geçişte ve 16 bitten 32 bite geçişte çok fark edilir bir hızlanma yaşanmıştır.
Ancak 32 bitten 64 bite geçişte hızlanma öncekiler kadar olmamıştır. Bunun nedeni 64 bitlik
işlemlerin aslında yoğun olarak yapılmadığındandır. Ancak 32 bitten 64 bite geçişin en önemli
etkisi işlemciye bağlanabilecek RAM miktarı üzerinde olmuştur. Örneğin elimizde 8 GB RAM’li
Raspberry Pi 5 olsun. Buradaki ARM cortex’i hem 32 bit hem de 64 bit işlemci gibi
çalışabilmektedir. Dolayısıyla biz bu kartımıza 32 bitlik de 64 bitlik de Linux yükleyebiliriz.
Ancak 32 bitlik Linux işletim sistemi işlemciyi 32 bit modda çalıştıracağı için işletim sistemi
8 GB RAM’in ancak 4 GB’sini kullanabilecektir. Bu 8 GB RAM’in tamamından istifade edebilmemiz
için işlemcinin 64 bit modunda çalıştırılması dolayısıyla da işletim sisteminin 64 bit olması
gerekecektir.

\sphinxAtStartPar
Peki 128 bitlik işlemciler tasarlansa, bunlar daha hızlı bir çalışma sunmaz mı? Neden 128 bitlik
işlemciler tasarlanmıyor? 128 bitlik tamsayı işlemlerine çok nadir gereksinim duyulmaktadır.
2$^{\text{128}}$ uzunluğunda RAM zaten şu an için erişilmesi imkânsız bir RAM miktarıdır. O halde
işlemcilerin 64 bitten 128 bite çıkartılmasının şu an için önemli bir faydası olmayacaktır.
Ancak tabii gelecekte böyle bir ihtiyaç ortaya çıkabilecektir.


\section{Popüler Ürünlerde ARM Cortex’leri}
\label{\detokenize{introduction:populer-urunlerde-arm-cortex-leri}}
\sphinxAtStartPar
Popüler ürünlerde kullanılan ARM cortex’leri şunlardır:
\begin{description}
\sphinxlineitem{\sphinxstylestrong{iPhone 14:}}
\sphinxAtStartPar
Apple’ın A15 isimli SoC’unu kullanmaktadır. Bu SoC’un içerisinde Apple’ın ARMv8.6\sphinxhyphen{}A
versiyonunu kullanan Avalanche Cortex’i vardır. Buradaki core’lar yalnızca 64 biti
desteklemektedir.

\sphinxlineitem{\sphinxstylestrong{iPhone 15:}}
\sphinxAtStartPar
Apple’ın A16 isimli SoC’unu kullanmaktadır. Bu SoC’un içerisinde Apple’ın ARMv8.6\sphinxhyphen{}A
versiyonunu kullanan Everest Cortex’i vardır. Buradaki core’lar yalnızca 64 biti
desteklemektedir.

\sphinxlineitem{\sphinxstylestrong{Apple M1 SoC:}}
\sphinxAtStartPar
Apple’ın ARMv8.4\sphinxhyphen{}A versiyonunu kullanan FireStorm isimli Cortex’i vardır. Buradaki core’lar
yalnızca 64 biti desteklemektedir.

\sphinxlineitem{\sphinxstylestrong{Apple M2 SoC:}}
\sphinxAtStartPar
Apple’ın ARMv8.6\sphinxhyphen{}A versiyonunu kullanan Avalanche ve Blizzard isimli Cortex’leri vardır.
Buradaki core’lar yalnızca 64 biti desteklemektedir.

\sphinxlineitem{\sphinxstylestrong{Apple M3 SoC:}}
\sphinxAtStartPar
Apple’ın ARMv8.6\sphinxhyphen{}A versiyonunu kullanan Avalanche ve Blizzard isimli Cortex’leri vardır.
Buradaki core’lar yalnızca 64 biti desteklemektedir.

\sphinxlineitem{\sphinxstylestrong{Samsung Galaxy S24:}}
\sphinxAtStartPar
Qualcomm firmasının Snapdragon\sphinxhyphen{}8 SoC’unu kullanmaktadır. Bu SoC içerisinde ARMv9\sphinxhyphen{}A versiyonunu
kullanan Cortex A715 bulunmaktadır.

\sphinxlineitem{\sphinxstylestrong{Redmi Note 11:}}
\sphinxAtStartPar
Xiaomi firmasının ürünüdür. Qualcomm Snapdragon\sphinxhyphen{}685 SoC’unu kullanmaktadır. Bu SoC’un
içerisinde ARMv8.2\sphinxhyphen{}A versiyonuna ilişkin Cortex A78 bulunmaktadır.

\end{description}


\section{Kursumuzda Kullanılan SBC’lerin ARM Cortex’leri}
\label{\detokenize{introduction:kursumuzda-kullanilan-sbc-lerin-arm-cortex-leri}}
\sphinxAtStartPar
Bu kursta kullanacağımız SBC’lerin üzerinde bulunan ARM cortex’leri ve bunların versiyon
numaraları şöyledir:

\begin{figure}[htbp]
\centering
\capstart

\noindent\sphinxincludegraphics[width=0.600\linewidth]{{sbc-arm-table}.png}
\caption{Kursumuzda kullanılan SBC’lerin ARM cortex karşılaştırması}\label{\detokenize{introduction:id5}}\end{figure}
\begin{description}
\sphinxlineitem{\sphinxstylestrong{Raspberry Pi 3:}}
\sphinxAtStartPar
BCM2837 SoC’u kullanılmıştır. Bu SoC’un içerisinde ARMv8\sphinxhyphen{}A versiyonuna ilişkin Cortex\sphinxhyphen{}A53
bulunmaktadır. Bu cortex’teki işlemciler hem 32 bit hem de 64 bit modda çalışabilmektedir.
Dolayısıyla bunlara 32 bit ve 64 bit Linux sistemleri yüklenebilmektedir.

\sphinxlineitem{\sphinxstylestrong{Raspberry Pi 4:}}
\sphinxAtStartPar
BCM2711 SoC kullanılmıştır. Bu SoC’un içerisinde ARMv8\sphinxhyphen{}A versiyonuna ilişkin Cortex\sphinxhyphen{}A72
bulunmaktadır. Bu cortex’teki işlemciler hem 32 bit hem de 64 bit modda çalışabilmektedir.
Dolayısıyla bunlara 32 bit ve 64 bit Linux sistemleri yüklenebilmektedir.

\sphinxlineitem{\sphinxstylestrong{Raspberry Pi 5:}}
\sphinxAtStartPar
BCM2712 SoC kullanılmıştır. Bu SoC’un içerisinde ARMv8.2\sphinxhyphen{}A versiyonuna ilişkin Cortex\sphinxhyphen{}A76
bulunmaktadır. Bu cortex’teki işlemciler hem 32 bit hem de 64 bit modda çalışabilmektedir.
Dolayısıyla bunlara 32 bit ve 64 bit Linux sistemleri yüklenebilmektedir.

\sphinxlineitem{\sphinxstylestrong{BeagleBone Black 4G ve BeagleBone Black Wireless:}}
\sphinxAtStartPar
Texas Instruments firmasının Sitara\sphinxhyphen{}AM3358 SoC’unu kullanmaktadır. Bu SoC’un içerisinde
Cortex\sphinxhyphen{}A8 kullanılmıştır. Bu cortex 32 bit işlemcilerden oluşmaktadır. Bunların kullandığı
komut mimarisi de ARMv7\sphinxhyphen{}A’dır. BeagleBone Black’ler 32 bit ARM işlemcilerini kullandığı için
bunlara yalnızca 32 bit Linux sistemleri yüklenebilmektedir.

\end{description}



\renewcommand{\indexname}{Dizin}
\printindex
\end{document}